\subsection{优化结果与动态辅助冷启动特性分析}
\label{sec:q4:results}

在上述预冷模型与闭环求解方法基础上，分别计算完全冷却、预冷20 min和预冷40 min三种工况。为兼顾辅助能耗与已测试扰动下的可行性，本节以留裕度动态方案\(D_R\)作为推荐策略，以沿用问题三功率分配并采用相同测量保持规则的恒功率方案\(C_H\)作为主要对照。两者使用相同的初始温度场、电流加载曲线和测量停机条件，比较结果见表\ref{tab:q4:main}。表中启动时间为真实状态首次满足成功判据的时刻\(t_s\)，辅助电能及电压、冰量和温差极值均统计至实际关热时刻\(t_{\rm off}\)；关热后的延续运行结果另行讨论。

\begin{table}[htbp]
    \centering
    \caption{不同预冷工况下动态与恒功率辅助启动结果对比}
    \vspace{0.3cm}
    \label{tab:q4:main}
    \renewcommand{\arraystretch}{1.5}
    \resizebox{\textwidth}{!}{
    \begin{tabular}{l c c c c c c c c c}
    \toprule
    \textbf{工况} & \textbf{策略} & \shortstack{\textbf{功率控制}\\\textbf{方式}} & \shortstack{\textbf{启动时间}\\\(t_s/\mathrm s\)} & \shortstack{\textbf{关热时刻}\\\(t_{\rm off}/\mathrm s\)} & \shortstack{\textbf{辅助电能}\\\(E_{\rm aux}/\mathrm J\)} & \shortstack{\textbf{最大温差}\\\(\Delta T_{\max}/\mathrm K\)} & \shortstack{\textbf{最低电压}\\\(V_{\min}/\mathrm V\)} & \shortstack{\textbf{最大局部}\\\textbf{冰体积分数}} & \textbf{结果}\\
    \midrule
    完全冷却 & \(C_H\) & 固定分片功率 & 25.741 & 28.2 & 3260.281 & 27.003 & 0.66185 & 0.012130 & 成功\\
    完全冷却 & \(D_R\) & 分片动态反馈 & 85.244 & 88.0 & 2929.996 & 6.060 & 0.56044 & 0.277253 & 成功\\
    预冷20 min & \(C_H\) & 固定分片功率 & 5.606 & 8.0 & 924.902 & 8.716 & 0.73621 & \(9.528\times10^{-5}\) & 成功\\
    预冷20 min & \(D_R\) & 分片动态反馈 & 74.308 & 78.2 & 293.609 & 9.054 & 0.60497 & 0.087293 & 成功\\
    预冷40 min & \(C_H\) & 固定分片功率 & 17.037 & 19.4 & 2242.888 & 19.852 & 0.68486 & 0.002818 & 成功\\
    预冷40 min & \(D_R\) & 分片动态反馈 & 88.539 & 91.6 & 1855.385 & 7.243 & 0.57332 & 0.239568 & 成功\\
    \bottomrule
    \end{tabular}}
\end{table}

由表\ref{tab:q4:main}可知，三种工况下两类策略均完成启动，但其能量投入和时间分配明显不同。动态策略将三工况辅助电能分别降至2929.996 J、293.609 J和1855.385 J，同时将启动过程延长至74--89 s，使反应热在升温过程中承担更大比例的热负荷。因此，动态策略的节能效果需要结合启动速度、片间一致性及结冰水平共同评价，不能仅根据辅助电能降低便认为全部性能同步改善。

\subsubsection{三种预冷工况的初始温度分布特征}

预冷结束时的温度场决定了后续启动所需补充的显热，同时影响各片电池与端板之间的初始传热方向。表\ref{tab:q4:initial}给出三种工况的代表节点温度和不同空间范围的温差。由于左右结构及边界近似对称，将两端对应节点合并列示；表中热容加权均温包含端板、双极板及MEA，全场温差包含端板真实外表面，五片温差则仅由五个单电池平均温度计算。

\begin{table}[htbp]
    \centering
    \caption{三种工况的初始温度及空间非均匀性}
    \vspace{0.3cm}
    \label{tab:q4:initial}
    \renewcommand{\arraystretch}{1.5}
    \resizebox{\textwidth}{!}{
    \begin{tabular}{l c c c c c c c}
    \toprule
    \textbf{工况} & \shortstack{\textbf{端板均温}\\\(/{}^\circ\mathrm C\)} & \shortstack{\textbf{第1、5片}\\\(/{}^\circ\mathrm C\)} & \shortstack{\textbf{第2、4片}\\\(/{}^\circ\mathrm C\)} & \shortstack{\textbf{第3片}\\\(/{}^\circ\mathrm C\)} & \shortstack{\textbf{热容加权均温}\\\(/{}^\circ\mathrm C\)} & \shortstack{\textbf{全场温差}\\\(/\mathrm K\)} & \shortstack{\textbf{五片温差}\\\(/\mathrm K\)}\\
    \midrule
    完全冷却 & -30.0000 & -30.0000 & -30.0000 & -30.0000 & -30.0000 & 0 & 0\\
    预冷20 min & -9.2554 & -9.0879 & -8.9790 & -8.9493 & -9.2123 & 0.50698 & 0.13860\\
    预冷40 min & -22.1591 & -22.0958 & -22.0547 & -22.0434 & -22.1428 & 0.19162 & 0.05239\\
    \bottomrule
    \end{tabular}}
\end{table}

\begin{figure}[!h]
\centering
\includegraphics[width=.95\textwidth]{图01_预冷温度场族.pdf}
\caption{不同预冷时间下的整堆温度场及局部温差}
\label{F7-1}
\end{figure}

由图\ref{F7-1}(a)及表\ref{tab:q4:initial}可见，预冷20 min和40 min时温度均呈中部较高、两端较低的分布，但其均温分别约为\(-9.21^\circ\mathrm C\)和\(-22.14^\circ\mathrm C\)，仍明显高于完全冷却工况。端板既承担对外散热，又占据整堆的大部分热容，因此控制了整体降温的主要时间尺度。离散热系统的最慢模态时间常数约为1233.26 s，即20.55 min，与总热容除以两端换热导度所得的集总估算接近。预冷从20 min延长至40 min后，电堆与环境之间的温差继续减小，整堆也逐步接近均匀的\(-30^\circ\mathrm C\)状态。

图\ref{F7-1}(b)将各时刻温度减去相应最低温度后显示局部梯度，可以看到MEA所在薄层的温度变化较集中，而高导热双极板内部温度变化较小。在附件给定物性和端面换热条件下，20 min的全场温差约为0.507 K，40 min时反而减小至0.192 K；对应的五片平均温差仅为0.139 K和0.052 K。因此，两种未完全冷却工况确实存在空间非均匀性，但不能由预冷时间延长直接推出绝对温差增大。本模型给出的主要变化是整体初温进一步降低，初始梯度在较长预冷阶段逐渐衰减。后续启动中若出现数开尔文乃至数十开尔文的片间温差，其主要来源应从端板吸热、分片加热和反应产热的动态竞争中解释。

\subsubsection{动态功率调节与温度、电压及结冰演化}

图\ref{F7-2}--图\ref{F7-4}分别给出三种工况推荐策略的功率、温度、电压和局部冰峰值。各图竖直分界线表示测量保持完成并锁存关热的时刻，阴影部分为关热后60 s的独立观察窗口。因此，评价表\ref{tab:q4:main}中的主结果时应使用分界线之前的数据，而不能把后验段的冰量增长或温度回落混入主窗口极值。

\begin{figure}[!h]
\centering
\includegraphics[width=.95\textwidth]{图04_case1_动态控制轨迹.pdf}
\caption{完全冷却工况下推荐动态策略的功率与水热电状态演化}
\label{F7-2}
\end{figure}

对于完全冷却工况，图\ref{F7-2}(a)显示，各片在启动初期均获得短时辅助热输入，随后内部第2--4片的功率逐渐降至零，而端部第1、5片的加热功率随热负荷变化持续增加，后期达到\(1\ \mathrm{W\,cm^{-2}}\)上限。内部电池能够利用自身反应热及片间热传递继续升温，端部电池则同时承担端板吸热与环境散热，因而成为决定全堆首次越过冰点的限制位置。动态控制通过减少内部电池的持续加热，将主要辅助电能集中于端部，使过程最大片间温差控制在6.060 K，显著低于固定恒功率方案的27.003 K。

\begin{figure}[!h]
\centering
\includegraphics[width=.95\textwidth]{图05_case2_动态控制轨迹.pdf}
\caption{预冷20 min工况下推荐动态策略的功率与水热电状态演化}
\label{F7-3}
\end{figure}

预冷20 min时，各部件仍保留较多显热，辅助加热需求明显减小。由图\ref{F7-3}(a)可见，端部功率先缓慢增加再逐渐回落，峰值仅为\(0.1243\ \mathrm{W\,cm^{-2}}\)，内部各片功率更低。端部电池累计实际加热时间为67.4 s，此后辅助功率已经降至零，而全堆首次达标时刻为74.308 s，测量保持确认时刻为78.2 s。这表明该工况末段主要依靠反应热和已形成的热状态完成升温，控制器无需在整个启动窗口持续供热。端部第1、5片在初期还存在短暂回落，说明较高初温虽然降低了总体热量缺口，但加载初期反应热较弱时，环境散热和端板吸热仍可能超过即时热输入。

\begin{figure}[!h]
\centering
\includegraphics[width=.95\textwidth]{图06_case3_动态控制轨迹.pdf}
\caption{预冷40 min工况下推荐动态策略的功率与水热电状态演化}
\label{F7-4}
\end{figure}

预冷40 min的初温介于前两种工况之间，图\ref{F7-4}(a)所示端部功率总体低于完全冷却工况，峰值约为\(0.5978\ \mathrm{W\,cm^{-2}}\)；内部第2--4片仅在早期加热约18.4 s，端部加热则维持至91.6 s。虽然该工况辅助电能低于完全冷却工况，但其首次启动时间88.539 s略长于完全冷却工况的85.244 s，原因是两种初态分别采用独立整定的控制参数，具有不同的功率水平和参考升温轨迹。因此，跨工况比较时，初温高低能够解释热量需求的变化，却不能脱离控制参数直接判断启动时间的排序。

为了量化空间功率分配，表\ref{tab:q4:cellenergy}给出推荐策略的逐片辅助电能。完全冷却、预冷20 min和预冷40 min工况中，端部两片合计分别占总辅助电能的约96.5\%、91.5\%和97.0\%。与问题三固定策略对第2、4片长期施加满功率不同，动态方案能够根据实际温升及时降低内部片功率，从而减少内部过热。这里的功率分配是片间热耦合与反馈共同作用的结果，并非预先规定只有端部允许加热。

\begin{table}[htbp]
    \centering
    \caption{推荐动态策略的逐片辅助电能分配}
    \vspace{0.3cm}
    \label{tab:q4:cellenergy}
    \renewcommand{\arraystretch}{1.5}
    \begin{tabular}{l c c c c c c}
    \toprule
    \textbf{工况} & \textbf{第1片/J} & \textbf{第2片/J} & \textbf{第3片/J} & \textbf{第4片/J} & \textbf{第5片/J} & \textbf{合计/J}\\
    \midrule
    完全冷却 & 1413.246 & 34.529 & 34.447 & 34.529 & 1413.246 & 2929.996\\
    预冷20 min & 134.347 & 9.674 & 5.566 & 9.674 & 134.348 & 293.609\\
    预冷40 min & 899.927 & 18.566 & 18.399 & 18.566 & 899.927 & 1855.385\\
    \bottomrule
    \end{tabular}
\end{table}

从图\ref{F7-2}--图\ref{F7-4}的电压和冰量子图可见，电流爬升过程中反应产水与低温冻结共同增加冰量，而温度升高和部分区域越过冰点后又促进融冰。推荐策略三工况的最低单片电压分别为0.56044 V、0.60497 V和0.57332 V，局部最大冰体积分数分别为0.277253、0.087293和0.239568，均满足0.30 V和0.99的约束。相比固定恒功率，动态策略在低温区停留更久、累计加载电荷更大，因此主窗口冰峰值更高，最低电压也更低。由此可见，本次优化是在保持安全可行的条件下减少辅助电能，并未同时实现最强的抑冰效果。

还需区分控制结构中设置的安全保护与实际轨迹中触发的动作。状态记录表明，三种推荐方案在无扰动主工况下均未进入安全增强状态，主要控制过程由温升跟踪、温升不足补偿及近目标降功率构成。安全增强提供的是风险上升时的后备规则，不能将其描述为这些主工况节能或抑冰的已发生原因。图中冰量采用全MEA各控制体的局部冰体积分数最大值，与MEA体积平均冰量及孔隙冰饱和度含义不同，后两者不能替代本节的安全约束指标。

\subsubsection{辅助能耗、启动速度与恒功率策略综合比较}

在相同测量保持规则下，定义推荐动态方案相对固定恒功率的辅助节电率为
\begin{equation}
\eta_E=\frac{E_{\rm aux}^{C_H}-E_{\rm aux}^{D_R}}{E_{\rm aux}^{C_H}}\times100\%.
\tag{7-@@R1@@}\label{eq:q4:saving}
\end{equation}
由表\ref{tab:q4:main}计算得到，完全冷却、预冷20 min和预冷40 min的辅助节电率分别为10.13\%、68.26\%和17.28\%。其中20 min工况初始热量储备较高，沿用完全冷却工况整定的固定大功率会产生较强的加热冗余，动态控制降低功率后的相对收益最大；完全冷却工况需补充的显热较多，端部功率后期达到上限，因此节电空间相对有限。

上述辅助节电伴随着明显的时间代价。固定恒功率三工况的首次启动时间分别为25.741 s、5.606 s和17.037 s，而推荐动态方案分别为85.244 s、74.308 s和88.539 s。推荐方案首次达标时的累计电荷依次为16.573、13.293和17.562\(\ \mathrm{C\,cm^{-2}}\)，虽然均未超过继承的20\(\ \mathrm{C\,cm^{-2}}\)预算，但明显高于固定恒功率的1.657、0.079和0.726\(\ \mathrm{C\,cm^{-2}}\)。更长的反应时间提供了更多反应热，同时也意味着消耗更多反应物。因此，式\eqref{eq:q4:saving}仅衡量电热丝辅助电能的变化，不能直接解释为氢耗或全系统总能耗的同幅下降。

片间一致性也存在工况差异。完全冷却和40 min工况的最大温差分别由27.003 K和19.852 K降至6.060 K和7.243 K，表明动态控制有效抑制了内部片过热；但20 min工况由8.716 K略增至9.054 K。较暖初态允许控制器长时间采用较小辅助功率，端部热损失与内部反应热的不均衡仍会逐步形成温差。因此，动态策略对温差的改善主要体现在本次两个较冷工况，不能概括为所有初态下一致性均优于恒功率。

\begin{figure}[!h]
\centering
\includegraphics[width=.95\textwidth]{图14_重优化恒功率与推荐动态.pdf}
\caption{固定恒功率、同工况重新优化恒功率与推荐动态策略的比较}
\label{F7-5}
\end{figure}

为进一步区分反馈调节本身的收益与原固定功率分配的适用性，图\ref{F7-5}补充了按各工况重新整定的恒功率方案\(C_*\)。在相同测量停机规则下，\(C_*\)三工况能耗分别为3237.876 J、0 J和1984.470 J，首次启动时间分别为29.416 s、84.725 s和95.287 s。推荐动态方案在完全冷却和40 min工况下相对\(C_*\)分别减少307.880 J和129.085 J，节电约9.51\%和6.50\%；其中40 min工况还将首次启动提前约6.75 s，但完全冷却工况仍需更长时间。

20 min工况的重新优化恒功率退化为五片零加热，名义动态优选方案也得到相同的0 J结果。这说明该初态在给定模型及加载条件下能够利用反应热实现自启动，辅助能耗的非负下界已经达到。推荐方案投入293.609 J，用于缩短启动时间并提高已测试扰动下的可行性，因此不能将其称为该工况的名义最小能耗，也不能将相对原固定功率的全部节电归因于反馈结构独有的优势。重新优化恒功率仅进行了名义参数搜索，其结果也不能代表经过同等扰动整定后的恒功率能力上限。

此外，本章固定功率的首次达标复算能耗约为2976.015 J，而采用测量保持规则后增至3260.281 J。前者与第六章约2975.3 J的结果接近，微小差异与数值设置及事件求解精度有关；后者增加的能耗主要来自温度裕度确认与保持期间继续加热。因此，本章节电率始终使用表\ref{tab:q4:main}中的共同停机基准，避免将第六章首次达标能耗与本章实际关热能耗直接相除。

\subsubsection{预冷时间对恒功率冷启动及单片一致性的影响}

针对问题四第（2）小问，将预冷时间依次取为10、15、\(\ldots\)、100 min，并始终施加问题三的固定功率分配和统一电流曲线。该组计算用于隔离预冷程度的影响，采用首次真实达标即结束的口径，不附加测量温度裕度和2 s保持。图\ref{F7-6}给出全部19个扫描点，表\ref{tab:q4:scan}选取其中六个代表时刻列出结果。

\begin{figure}[!h]
\centering
\includegraphics[width=.95\textwidth]{图08_预冷时间扫描.pdf}
\caption{预冷时间对固定恒功率启动性能的影响}
\label{F7-6}
\end{figure}

\begin{table}[htbp]
    \centering
    \caption{不同预冷时间下固定恒功率启动的代表性结果}
    \vspace{0.3cm}
    \label{tab:q4:scan}
    \renewcommand{\arraystretch}{1.5}
    \resizebox{\textwidth}{!}{
    \begin{tabular}{c c c c c c c c}
    \toprule
    \shortstack{\textbf{预冷时间}\\\(/\mathrm{min}\)} & \shortstack{\textbf{热容加权初温}\\\(/{}^\circ\mathrm C\)} & \shortstack{\textbf{首次达标}\\\(/\mathrm s\)} & \shortstack{\textbf{辅助电能}\\\(/\mathrm J\)} & \shortstack{\textbf{最大五片温差}\\\(/\mathrm K\)} & \shortstack{\textbf{最低电压}\\\(/\mathrm V\)} & \shortstack{\textbf{最大局部}\\\textbf{冰体积分数}} & \textbf{结果}\\
    \midrule
    10 & 3.812 & 0 & 0 & 0.225 & 1.23732 & 0 & 初态已达标\\
    20 & -9.212 & 5.606 & 648.136 & 6.298 & 0.74738 & \(9.528\times10^{-5}\) & 成功\\
    40 & -22.143 & 17.037 & 1969.681 & 17.469 & 0.68684 & 0.002818 & 成功\\
    60 & -27.030 & 22.316 & 2579.978 & 22.416 & 0.67089 & 0.007411 & 成功\\
    80 & -28.878 & 24.428 & 2824.169 & 24.161 & 0.66534 & 0.010141 & 成功\\
    100 & -29.576 & 25.242 & 2918.317 & 24.799 & 0.66329 & 0.011347 & 成功\\
    \bottomrule
    \end{tabular}}
\end{table}

由图\ref{F7-6}(a)--(c)可见，预冷越充分，整堆初温越接近\(-30^\circ\mathrm C\)，固定功率启动所需时间和辅助电能也随之增加，但增加幅度逐渐减小。10 min时五片初温约为4.01--4.24\(^\circ\mathrm C\)，起始状态已经满足温度条件，因此首次达标时间和辅助能耗均为零；该点不代表经历了一个零时间的加热过程，也不构成任意后续负载下可持续运行的证明。20 min和40 min的辅助电能分别为648.136 J和1969.681 J，而100 min时达到2918.317 J，逐步接近完全冷却工况约2976.015 J的首次达标基准。由于所有扫描点的总加热功率保持不变，能耗增长与首次启动时间延长直接对应，其根源是预冷后需要补充的显热增加。

图\ref{F7-1}与图\ref{F7-6}(d)还揭示了初始温差和启动温差的不同变化方向。随着预冷时间由10 min增加到100 min，完整预冷场的绝对温差由约0.825 K减至0.010 K，五片初始均温差由约0.225 K减至0.003 K；但启动过程最大片间温差却由0.225 K增加至24.799 K。较深预冷使大热容端板处于更低温度，启动时端部电池需要持续向端板供热，而内部电池在固定加热和反应产热共同作用下较快升温。因而，较均匀的低温初态并不意味着较均匀的启动过程，启动期间重新形成的端部热负荷差异才是后期片间不一致性的主要原因。

在电压与冰量方面，预冷时间增加使低温产水持续时间延长，局部冰峰值逐渐升高，最低单片电压逐渐降低。即使预冷100 min，最低电压仍为0.66329 V，局部冰峰值仅为0.011347，均保持在规定阈值内。全部19个离散预冷时间点均完成首次启动，说明问题三固定功率方案在所考察区间内具有可行性；本次扫描没有发现失败点，因此不能据此给出未经识别的临界预冷时间。该结论仅针对已计算的19个点及固定参数条件。

\subsubsection{关热后温度回落与持续运行特性}

首次越过冰点并完成短时保持，只说明启动阶段满足规定判据，并不等价于整个电堆已经具备维持暖态的热储备。为检验这一差异，在各方案实际关热后继续按题给电流曲线运行60 s，辅助功率保持锁存为零。表\ref{tab:q4:post}给出推荐方案关热时的状态及后验段极值，同时列出固定恒功率方案的后验最低温度作为参照。

\begin{table}[htbp]
    \centering
    \caption{推荐动态策略关热状态及关热后60 s的运行结果}
    \vspace{0.3cm}
    \label{tab:q4:post}
    \renewcommand{\arraystretch}{1.5}
    \resizebox{\textwidth}{!}{
    \begin{tabular}{l c c c c c c}
    \toprule
    \textbf{工况} & \shortstack{\textbf{关热最低片温}\\\(/{}^\circ\mathrm C\)} & \shortstack{\textbf{关热左端板温度}\\\(/{}^\circ\mathrm C\)} & \shortstack{\textbf{后验最低片温}\\\(/{}^\circ\mathrm C\)} & \shortstack{\textbf{后验最低电压}\\\(/\mathrm V\)} & \shortstack{\textbf{后验局部}\\\textbf{冰峰值}} & \shortstack{\(C_H\)\textbf{后验最低片温}\\\(/{}^\circ\mathrm C\)}\\
    \midrule
    完全冷却 & 0.882 & -14.471 & -4.508 & 0.55967 & 0.50104 & -7.986\\
    预冷20 min & 0.433 & -5.115 & 0.433 & 0.61376 & 0.08702 & -2.677\\
    预冷40 min & 0.763 & -10.288 & -1.910 & 0.58853 & 0.29566 & -6.556\\
    \bottomrule
    \end{tabular}}
\end{table}

由表\ref{tab:q4:post}及图\ref{F7-2}、图\ref{F7-4}的阴影部分可见，完全冷却和40 min工况虽然在关热时全部单电池均已高于冰点，但端板仍分别约为\(-14.47^\circ\mathrm C\)和\(-10.29^\circ\mathrm C\)。辅助热源一旦撤去，端板继续从端部单电池吸热，使最低片温分别降至\(-4.508^\circ\mathrm C\)和\(-1.910^\circ\mathrm C\)。这与第六章发现的端板热惯性效应一致，说明仅根据单电池温度过零及2 s保持，尚不能判定端板热汇已经充分减弱。

相比之下，20 min工况的初始热储备更高，且在正式关热之前辅助功率已提前降至零，反应热能够维持其继续升温；关热后最低片温保持在0.433\(^\circ\mathrm C\)以上。完全冷却和40 min工况虽再次进入零下区间，但后验最低电压仍高于0.30 V，冰峰值也未达到0.99，因此应表述为未维持持续暖态，而不是将其与电压失效或严重冰堵混为一谈。与固定恒功率方案相比，推荐动态方案在本次60 s观察窗口内减轻了回冷，但仍未在全部工况消除该现象。

上述后验段不计入主表辅助电能，且不允许重新开启加热。若进一步要求关热后始终保持冰点以上，则需将这一要求明确纳入优化约束，并重新整定功率与热裕度；本次结果不能直接外推为已满足该更强目标。后验段继续加载也会消耗反应物，其作用是检验延续运行状态，不能混入首次启动的电荷预算。

\subsubsection{参数灵敏度与控制策略稳健性分析}

图\ref{F7-7}分别展示预冷传热参数变化和名义动态控制参数扰动的影响。预冷部分同时考察绝对温差及其相对于整堆平均超温的比值，以区分内部传热不均匀程度与整体冷却速度；启动部分则固定名义控制结构，对片间导度、端板耦合导度、环境换热及控制参数逐项扰动。

\begin{figure}[!h]
\centering
\includegraphics[width=.95\textwidth]{图11_物理参数与控制器敏感性.pdf}
\caption{预冷传热参数与名义动态控制参数的敏感性}
\label{F7-7}
\end{figure}

由图\ref{F7-7}(a)、(b)可见，内部导热能力降低时，端部与中部更难迅速交换热量，预冷温度场的相对非均匀程度增大。端面换热增强则同时提高外部散热强度并加快整体降温，因此固定预冷时刻的绝对温差不一定随换热系数单调增加；当电堆已经较接近环境温度时，可用于维持空间梯度的超温也随之减小。采用温差与热容加权均温超出环境温度之比，可更清楚地识别外部换热与内部导热竞争的影响。

启动敏感性表明，端板耦合导度对辅助电能的影响较突出。以完全冷却的名义方案为例，将端板耦合导度降低20\%后，辅助电能约降至2453.251 J；提高20\%后则增至约3190.957 J，显著偏离基准2855.408 J。端板耦合越强，启动过程中单电池向低温端板传递的热量越多，维持参考温升所需的辅助电能也越高。另一方面，个别控制参数扰动会导致未能完成测量保持停机，故图中的失败候选即使已消耗较少电能，也不能作为更优的成功方案。图\ref{F7-7}使用名义动态参数，其作用是揭示影响机制，不代表推荐方案已经在全部参数组合中得到严格鲁棒保证。

\begin{figure}[!h]
\centering
\includegraphics[width=.95\textwidth]{图13_名义与留裕度备选.pdf}
\caption{名义节能方案与推荐留裕度方案的能耗、时间及扰动表现}
\label{F7-8}
\end{figure}

由图\ref{F7-8}(a)、(b)可见，名义方案\(D_N\)三工况辅助电能为2855.408 J、0 J和1767.146 J，其中两个较冷工况的实际关热时刻均达到98.4 s，靠近允许的保持窗口上限。推荐方案\(D_R\)分别增加约74.588 J、293.609 J和88.239 J，将关热时刻提前至88.0 s、78.2 s和91.6 s。其目的在于为初温偏差、滤波滞后及物性失配保留可用余量，而不是在名义模型下进一步压低辅助能耗。

名义方案在原测量噪声与初场偏差试验中，仅20/90次满足全部可行性条件，成功样本集中于20 min工况；推荐方案在参数冻结后的独立噪声与初场试验中通过90/90次，并在24个物性失配与噪声联合场景中全部通过。两组噪声试验采用不同随机种子，因而这些比例可以说明各方案在所列场景下的表现，但不能作为严格配对的统计优势证明。固定恒功率\(C_H\)在相同最终验证场景下也全部通过，说明它本身具有较充分的启动时间余量。推荐动态方案的价值在于在这些已验证场景中保留可行性的同时降低辅助电能，而不是声称只有动态控制才能抵御扰动。

\begin{figure}[!h]
\centering
\includegraphics[width=.95\textwidth]{图16_观测器误差与测量停机.pdf}
\caption{参数失配与测量噪声下的观测误差及停机事件辨析}
\label{F7-9}
\end{figure}

图\ref{F7-9}进一步检验反馈信息的可获得性。图中采用未参与训练的参数失配与测量噪声样例，控制器仅使用独立观测模型和温度、电压采样。端板温度估计随物性失配出现偏差，说明观测状态并未直接复制被控对象的真实状态；冰风险估计与真实局部冰峰值之间也保留误差，不能视为对完整冰场的精确重构。以图\ref{F7-9}(c)的完全冷却样例为例，真实首次达标约为90.70 s，而滤波测温满足温度裕度并连续保持后，到95.0 s才确认关热；20 min样例的相应时刻约为82.61 s和91.0 s。这种差异来自传感采样、滤波和保持计时，说明实际关热时刻不能简单设为\(t_s+2\ \mathrm s\)。

因此，推荐方案的稳健性应由真实状态下的成功时刻、路径约束和关热前真实状态连续保持是否完成共同判定，不能仅以测量判据被触发为依据。图\ref{F7-8}(d)仍显示部分工况关热后回冷，也再次说明本节的扰动可行性验证与持续暖态验证属于不同层次。

\subsubsection{能量守恒、数值收敛与结果可靠性分析}

为解释节电的能量来源，将启动主窗口中的辅助加热、反应产热、相变净释热、环境散热和离散显热增量分别累计。按照与时间推进相同的热容和时间积分定义，离散能量残差写为
\begin{equation}
R_E=E_{\rm aux}+E_{\rm gen}+E_{\rm phase}-E_{\rm loss}-E_{\rm sensible}.
\tag{7-@@R2@@}\label{eq:q4:energyresidual}
\end{equation}
其中\(E_{\rm phase}\)为带符号的相变净热，\(E_{\rm sensible}\)按各时间步使用的有效热容与温度增量累计。该量与离散方程保持一致，不将随含水状态变化的有效热容简单替换为一个常数乘总温升。

\begin{figure}[!h]
\centering
\includegraphics[width=.95\textwidth]{图09_能量收支与累积能耗.pdf}
\caption{推荐动态与固定恒功率策略的能量收支及累计辅助电能}
\label{F7-10}
\end{figure}

由图\ref{F7-10}(a)可见，各工况热输入总量与显热增量、环境散热之和相互吻合。完全冷却推荐方案在关热前输入辅助热2929.996 J、反应热1941.572 J和相变净热13.454 J，扣除环境散热272.961 J后，离散显热增量为4612.062 J。相同工况固定恒功率方案的反应热仅为194.959 J，而辅助输入为3260.281 J。因此，动态方案虽然延长了升温过程并增加环境散热，却通过更充分地利用反应热降低了所需辅助电能，同时使关热时端板温度较高。这也解释了其关热后回冷较固定策略有所减轻。

对于预冷20 min工况，推荐方案的反应热达到1533.862 J，明显高于293.609 J的辅助热；相变净热为\(-2.020\) J，环境散热为394.056 J，最终显热增量约为1431.395 J。预冷40 min时，辅助热和反应热分别为1855.385 J与2027.484 J，二者共同提供主要升温能量。图\ref{F7-10}(b)中，20 min工况累计辅助电能在末段转为水平，而其余两种工况仍随端部加热持续增加，与前述功率轨迹及提前停止实际加热的现象一致。

三工况推荐策略的离散能量残差绝对值约为\(10^{-11}\)--\(10^{-10}\) J，完整轨迹水量残差最大绝对值约为\(10^{-16}\)--\(10^{-15}\) kg，表明当前离散更新中的收支能够闭合。预冷阶段另行检查全堆内能减少与两端累计散热的一致性，并利用同网格矩阵指数解校验时间离散。守恒残差很小主要说明数值收支自洽，仍需结合独立参考和网格加密判断离散误差，不能单凭该残差推断物理预测没有误差。

\begin{figure}[!h]
\centering
\includegraphics[width=.95\textwidth]{图10_网格步长与控制周期检验.pdf}
\caption{预冷温度场与推荐动态启动结果的时间和空间收敛检验}
\label{F7-11}
\end{figure}

由图\ref{F7-11}(a)、(b)可见，预冷时间步从2 s减小至0.25 s时，20 min温度场相对同网格矩阵指数解的最大误差由约0.0166 K降至0.00208 K，继续减小至0.125 s后约为0.00104 K，表现出后向欧拉方法预期的收敛趋势。增加空间控制体数量后，七节点投影相对加密参考的误差也进一步下降，说明预冷初始场的离散精度能够满足三工况差异分析的需要。

对启动过程，固定正式控制参数和0.2 s采样周期，将每片水相控制体由58个加密至232个，同时把时间步从0.025 s减小至0.00625 s。三工况辅助电能分别由2929.996 J、293.609 J和1855.385 J变为2928.508 J、293.469 J和1854.436 J，相对变化均约为0.05\%；首次启动时间变化均小于0.008 s，实际关热时刻保持不变。图\ref{F7-11}(c)由此表明，主要能耗和时间结论对进一步联合加密较稳定。

图\ref{F7-11}(d)同时显示，局部冰峰值比总能耗更依赖空间分辨率。完全冷却工况由0.27725升至0.28658，20 min工况由0.08729升至0.08912，40 min工况由0.23957升至0.24623。尽管加密后的冰量仍远低于0.99约束，且不改变三工况的可行性判断，但其相对变化大于辅助电能的变化，因此不能将所有局部状态均描述为完全网格无关。综合基准复现、守恒检查、独立参考、离散加密与扰动验证，本章结果可支持给定模型及搜索范围内的性能比较；预冷和启动阶段沿用的有效物性、换热位置近似，以及尚缺新增实验校准的限制仍需保留。

\subsection{问题四小结}

针对不同预冷程度下电堆动态辅助启动问题，本文在前述水--冰--热--电耦合模型基础上，建立了整堆预冷传热模型，将预冷温度场映射为五片电池与两端端板的非均匀初态；进一步利用温度、电压及其变化率构造状态识别与独立观测器，通过前馈、PI反馈、温升不足补偿、安全增强和测量保持停机形成分片动态加热策略。以辅助电能为首要目标完成有限参数搜索后，再根据扰动训练场景选择留有温度和时间裕度的推荐方案。

结果表明，推荐方案在完全冷却、预冷20 min和预冷40 min工况下分别消耗2929.996 J、293.609 J和1855.385 J辅助电能，相对共同测量停机规则下的问题三固定功率方案节电10.13\%、68.26\%和17.28\%。其节能主要来自减少内部片持续加热，并利用较长启动过程中的反应热分担升温需求，代价是启动较慢且主窗口结冰水平较高；三工况电压、冰量及启动电荷仍满足约束。推荐方案通过所列独立扰动场景，但完全冷却与40 min工况关热后仍出现回冷，因此完成首次启动及短时保持并不等同于实现持续暖态。

在10--100 min预冷扫描中，19个离散工况均完成固定恒功率首次启动。随着预冷加深，初温降低，启动时间、辅助电能和过程最大片间温差增加，并趋近完全冷却结果；预冷场自身的绝对温差则逐渐减小。这说明初始储热量决定总体升温需求，而启动期间的片间一致性还受到端板热汇和分片热输入的影响。
